\documentclass[12pt]{article}
\usepackage[utf8]{inputenc}
\usepackage[T1]{fontenc}
\usepackage[margin=1in]{geometry}
\usepackage{setspace}
\usepackage{xcolor}
\usepackage[colorlinks=true, urlcolor=blue!60!black, linkcolor=blue!60!black]{hyperref}
\setlength{\parindent}{0em}
\setlength{\parskip}{0.6em}
\renewcommand{\baselinestretch}{1.0}
\begin{document}
\begin{center}
{\LARGE \textit{South-Central Timber Development, Inc.\ v.\ Wunnicke}, 467 U.S.\ 82 (1984)\par}
\vspace{0.5em}
{\large Case Brief --- Two-Page Edition\par}
\end{center}

\section{Memory Jogger}
A condition requiring buyers of state timber to process it in Alaska was neither authorized by Congress nor shielded as market participation (\href{https://supreme.justia.com/cases/federal/us/467/82/}{Justia}).

\section{Detailed Case Facts}
Alaska sold state timber on condition that buyers process it in Alaska before export, accepting lower bids in exchange. South-Central Timber, which exported unprocessed logs to Japan, challenged the condition. Federal rules imposed a similar requirement on timber from federal lands.

\section{Procedural History}
The district court enjoined the condition; the Ninth Circuit reversed, finding implicit congressional authorization.

\section{Judicial Votes}
Reversed, 6--2. White wrote for the Court on authorization (with Burger, Brennan, Blackmun, Powell, and Stevens) and for a plurality of four on the rest; Powell and Burger concurred in part; Rehnquist and O'Connor dissented; Marshall did not participate.

\section{Holding}
Congressional consent to state trade barriers must be ``unmistakably clear,'' and a federal-lands policy is not. The plurality held that a market participant may not impose conditions regulating a downstream market.

\section{Analysis of Opinions}
\textbf{Majority:} Requiring a clear national decision protects out-of-state interests against Balkanization. \textbf{Plurality:} unlike \textit{Alexandria Scrap}, \textit{Reeves}, and \textit{White}, Alaska controlled buyers after the sale in the processing market, where it did not participate; the local-processing mandate is virtually \textit{per se} invalid, especially in foreign commerce. \textbf{Rehnquist:} the line is ``both artificial and unconvincing,'' since Alaska simply paid buyers through a lower price to hire local processors.

\section{Examples: Future Applications}
\textbf{Same-Side:} state-land lithium leases requiring in-state refining. \textbf{Opposite-Side:} a state-owned refinery selling its output only to local buyers. \textbf{Fence-Sitter:} requiring resident crews to log state land.

\section{Critique}
The plurality defined the market by ``intuition,'' and only four Justices joined it; but the rule pushes States toward visible subsidies.

\section{Key Quotations}
\begin{quote}
``There is no talismanic significance to the phrase `expressly stated,' however; it merely states one way of meeting the requirement that, for a state regulation to be removed from the reach of the dormant Commerce Clause, congressional intent must be unmistakably clear.''
\end{quote}
\noindent --- Justice White, for the Court, 467 U.S.\ at 91--92.

\end{document}
